等几何分析采用统一样条基函数同时描述 CAD 几何与分析场，为结构尺寸、形状和拓扑优化提供了几何建模、数值分析与灵敏度计算一体化框架。面向复杂壳体、曲边胞元和增材制造轻量化构件，传统“几何重建-网格再生-优化更新”流程往往存在模型转换频繁、边界质量受限和设计变量语义不一致等问题；等几何分析能够在统一几何表示下组织分析求解、优化更新和几何回写，为高保真结构优化提供技术基础。本文聚焦智能结构优化设计场景，综述等几何分析在结构优化中的主要进展。首先介绍 B 样条、NURBS 与等几何离散的基本表达，并概述 SIMP 材料插值、伴随灵敏度、滤波正则化、水平集演化和自由形变映射等经典模型；随后从密度法、水平集法、移动可变形构件和多补丁壳结构优化等路线梳理等几何拓扑优化与形状优化的代表性工作；最后讨论超材料、点阵结构、增材制造、智能生成和开源工作流等融合趋势。现有研究表明，复杂多补丁几何、局部细化、可制造性约束、几何回写以及数据驱动方法与物理模型的耦合，仍是推动该方向工程应用的关键问题。
